\subsection{Regular rings and flat extensions}

PROPOSITION 8.--- Let \(\rho: A \to B\) a homomorphism of Noetherian rings making B into a faithfully flat A-module.

\begin{enumerate}
\def\labelenumi{\alph{enumi})}
\item
  For every \(\Lambda\) -module M of finite type, we have \(\mathrm{dp}_{\mathrm{A}}(\mathrm{M}) = \mathrm{dp}_{\mathrm{B}}(\mathrm{B} \otimes_{\mathrm{A}} \mathrm{M})\) .
\item
  If the ring B is regular, the ring A is regular.
\end{enumerate}

The B-module \(B \otimes_{A} M\) is of finite type (I, §3, no. 6, prop. 12). For it to be zero, it is necessary and sufficient that M be so (I, §3, no. 1, def. 1); for it to be projective, it is necessary and sufficient that M be a projective A-module (I, no. 6, prop. 12). This proves a) when \(\mathrm{dp}_{\Lambda}(\mathrm{M}) \leqslant 0\) . Suppose therefore \(\mathrm{dp}_{\mathrm{A}}(\mathrm{M}) \geqslant 1\) (whence \(\mathrm{dp}_{\mathrm{B}}(\mathrm{B} \otimes_{\Lambda} \mathrm{M}) \geqslant 1\) ) and let us prove a) by induction on \(\mathrm{dp}_{\mathrm{A}}(\mathrm{M})\) . Choose an exact sequence of A-modules

\[
0 \to \mathrm{N} \to \mathrm{L} \to \mathrm{M} \to 0
\]

where the A-module L is free of finite type. We have \(\mathrm{dp}_{\Lambda}(\mathbf{N}) = \mathrm{dp}_{\mathbf{A}}(\mathbf{M}) - 1\) As B is flat over A, the sequence

\[
0 \rightarrow \mathrm{B} \otimes_ {\mathrm{A}} \mathrm{N} \rightarrow \mathrm{B} \otimes_ {\mathrm{A}} \mathrm{L} \rightarrow \mathrm{B} \otimes_ {\mathrm{A}} \mathrm{M} \rightarrow 0
\]

is exact and we have \(\mathrm{dp}_{\mathrm{B}}(\mathrm{B}\otimes_{\mathrm{A}}\mathrm{N}) = \mathrm{dp}_{\mathrm{B}}(\mathrm{B}\otimes_{\mathrm{A}}\mathrm{M}) - 1\) The induction hypothesis applied to N allows us to conclude. This proves a); assertion b) follows from a) and prop. 4 of no. 2.

COROLLARY. Let B be a regular integral domain and let A be a Noetherian subring of B such that B is a finitely generated A-module. The following conditions are equivalent:

\begin{enumerate}
\def\labelenumi{(\roman{enumi})}
\item
  A is regular;
\item
  B is a projective Λ-module;
\item
  B is a flat A-module;
\item
  B is a faithfully flat A-module.
\item
  \(\Rightarrow\) (ii): this follows from the corollary of Proposition 5 in No. 3.
\item
  \(\Rightarrow\) (iii): this follows from I, § 3, no. 1, prop. 1.
\item
  \(\Rightarrow\) (iv): for every prime ideal \(\mathfrak{p}\) of A, we have \(\mathfrak{pB} \neq B\) (V, § 2, no. 1, cor. 1 of Theorem 1). It then suffices to apply I, § 3; no. 1, prop. 1.
\item
  \(\Rightarrow\) (i): this follows from prop. 8, b).
\end{enumerate}

For every Noetherian local ring \( A \), let \( \delta(A) \) denote the integer

\[
\delta (\mathrm{A}) = \left[ \mathfrak {m} _ {\mathrm{A}} / \mathfrak {m} _ {\mathrm{A}} ^ {2}: \kappa_ {\mathrm{A}} \right] - \dim (\mathrm{A}).
\]

Recall (VIII, § 5, no. 1) that \(\delta(A)\) is always positive, and its vanishing characterizes regular local rings.

Let \(\rho : A \to B\) a local homomorphism of Noetherian local rings; we deduce from \(\rho\) be a \(\kappa_{A}\) -linear \(\mathfrak{m}_{A}/\mathfrak{m}_{A}^{2} \longrightarrow \mathfrak{m}_{B}/\mathfrak{m}_{B}^{2}\) , whence a homomorphism \(\kappa_{B}\) -linear

\[
d \rho : \kappa_ {B} \otimes_ {\kappa_ {A}} \left(\mathfrak {m} _ {A} / \mathfrak {m} _ {A} ^ {2}\right) \longrightarrow \mathfrak {m} _ {B} / \mathfrak {m} _ {B} ^ {2}.
\]

Lemma 1.--- We have

\[
\delta (\mathrm{B}) + \left[ \operatorname{Ker} (d \rho): \kappa_ {\mathrm{B}} \right] = \delta (\mathrm{A}) + \delta \left(\kappa_ {\mathrm{A}} \otimes_ {\mathrm{A}} ^ {\prime} \mathrm{B}\right) + (\dim (\mathrm{A}) - \dim (\mathrm{B}) + \dim \left(\kappa_ {\mathrm{A}} \otimes_ {\mathrm{A}} \mathrm{B}\right)).
\]

Let C denote the local ring \(\kappa_{\mathrm{A}}\otimes_{\mathrm{A}}\mathrm{B}\) . Consider the exact sequence of B-modules

\[
\mathrm{B} \otimes_ {\mathrm{A}} \mathfrak {m} _ {\mathrm{A}} \longrightarrow \mathfrak {m} _ {\mathrm{B}} \longrightarrow \mathfrak {m} _ {\mathrm{C}} \rightarrow 0;
\]

by tensor product with \(\kappa_{\mathrm{B}}\) one obtains an exact sequence of \(\kappa_{\mathrm{B}}\) -vector spaces

\[
\kappa_ {\mathrm{B}} \otimes_ {\kappa_ {\Lambda}} (\mathfrak {m} _ {\mathrm{A}} / \mathfrak {m} _ {\mathrm{A}} ^ {2}) \stackrel {d \rho} {\longrightarrow} \mathfrak {m} _ {\mathrm{B}} / \mathfrak {m} _ {\mathrm{B}} ^ {2} \longrightarrow \mathfrak {m} _ {\mathrm{C}} / \mathfrak {m} _ {\mathrm{C}} ^ {2} \rightarrow 0,
\]

whence one deduces the equality

\[
\left[ \mathfrak {m} _ {\mathrm{B}} / \mathfrak {m} _ {\mathrm{B}} ^ {2}: \kappa_ {\mathrm{B}} \right] + \left[ \operatorname{Ker} (d \boldsymbol {\rho}): \kappa_ {\mathrm{B}} \right] = \left[ \mathfrak {m} _ {\mathrm{A}} / \mathfrak {m} _ {\mathrm{A}} ^ {2}: \kappa_ {\mathrm{A}} \right] + \left[ \mathfrak {m} _ {\mathrm{C}} / \mathfrak {m} _ {\mathrm{C}} ^ {2}: \kappa_ {\mathrm{C}} \right],
\]

which implies the lemma.

PROPOSITION 9.- Let \(\rho : A \to B\) be a local homomorphism of Noetherian local rings. The following conditions are equivalent:

\begin{enumerate}
\def\labelenumi{(\roman{enumi})}
\tightlist
\item
  the ring B is regular and the map \(\kappa_{\mathrm{B}}\) -linear

\[
d \rho : \kappa_ {B} \otimes_ {\kappa_ {A}} \left(\mathfrak {m} _ {A} / \mathfrak {m} _ {A} ^ {2}\right) \longrightarrow \mathfrak {m} _ {B} / \mathfrak {m} _ {B} ^ {2}
\]

is injective;

\item
  the rings B and \(\kappa_{\mathrm{A}}\otimes_{\mathrm{A}}\mathrm{B}\) are regular and the A-module B is flat;
\item
  the rings A and \(\kappa_{\mathrm{A}}\otimes_{\mathrm{A}}\mathrm{B}\) are regular and the A-module B is flat;
\item
  the rings A and \(\kappa_{\Lambda} \otimes_{\mathrm{A}} \mathrm{B}\) are regular, and one has

\[
\dim (\mathrm{B}) = \dim (\mathrm{A}) + \dim (\kappa_ {\mathrm{A}} \otimes_ {\Lambda} \mathrm{B}).
\]
\end{enumerate}

We have \(\dim(B) \leqslant \dim(A) + \dim(\kappa_A \otimes_A B)\) (VIII, § 3, n° 4, cor. 1 of prop. 7); the equivalence of (i) and (iv) therefore follows from Lemma 1. Under hypothesis (ii), the A-module B is faithfully flat since we have \(\mathfrak{m}_A B \subset \mathfrak{m}_B \neq B\) (I, § 3, no. 1, prop. 1), which implies (iii) by prop. 8. The implication (iii)⇒(iv) follows from VIII, loc. cit.

It now suffices for us to prove that when the equivalent conditions (i) and (iv) are satisfied, the A-module B is flat. Let x be a system of coordinates of A. Since \(d\rho\) is injective, the image under \(\rho\) of this sequence is part of a coordinate system of B. Thus the sequence x is completely secant for \(\Lambda\) and for B (VIII, § 5, no. 3, prop. 2), and generates the ideal \(\mathfrak{m}_{\Lambda}\) from A. According to remark 3 of A, X, p.~159, the A-module \(\mathrm{Tor}_{1}^{\mathrm{A}}(\kappa_{\mathrm{A}},\mathrm{B})\) is isomorphic to \(\mathrm{H}_{1}(\mathbf{x},\mathrm{B})\) ; hence it is zero; it follows that B is flat over A (III, § 5, no. 2, th. 1 and no. 4, prop. 2).

Example.--- * Let X, Y be two complex analytic varieties, locally of finite dimension, \(f\) a morphism from X to Y, and \(x\) a point of X. Consider the local homomorphism \(\rho: \mathcal{O}_{\mathrm{Y}, f(x)} \to \mathcal{O}_{\mathrm{X}, x}\) associated with \(f\) . The map \(d\rho\) is the transpose of the tangent map \(\mathrm{T}_x(f): \mathrm{T}_x(\mathrm{X}) \to \mathrm{T}_{f(x)}(\mathrm{Y})\) The conditions (i) to (iv) of Proposition 9 are therefore equivalent in this case to the fact that \(f\) is a submersion at \(x\) (VAR, R, 5.9.1).*

COROLLARY.--- Let \(\rho: \Lambda \to B\) a homomorphism of Noetherian rings making B a flat A-module. If A is regular and if \(\kappa(\rho^{-1}(\mathfrak{n})) \otimes_{\Lambda} B\) is regular for every maximal ideal \(\mathfrak{n}\) of B, the ring B is regular.

Indeed for every maximal ideal \(\mathfrak{n}\) of B the \(A_{\rho^{-1}(n)}\) -module \(B_{n}\) is flat (II, § 3, no. 4, prop. 15), so that the ring \(B_{n}\) is regular by Prop. 9.

